\section{Conventions and elementary corrections in the inputs}
\label{sec:input-conventions}

The proof uses the versioned statements of \cite{DKLM2026} specified
in the text. We record elementary corrections needed when reading their
proofs literally, and give the conformal estimate used in
\cref{lem:flat-patch}. These are our reconciliations of the indicated
formulas, not an author-issued erratum. They do not replace the
substantive regularity, convergence, or free-energy theorems.

\subsection{The path decomposition in the increment estimate}
\label{sec:input-path}

In \cite[Lemma~21.3]{DKLM2026}, the printed path-length factor $S'$
must be $\max\{1,-S'\}$, as used in its subsequent Equation~(367).
Indeed $S'$ may be negative, whereas a nonconstant path has positive
length. Here is a direct construction of the required path.

Order the increments so that $|a-a'|\le|b-b'|$. If $a=a'$, use
the empty path. Otherwise put $R=|a-a'|>0$, $B=\{b,b'\}$,
$d=\dist(\{a,a'\},B)$, and
$D=\max\{1,d\}$. Fix an order $k$ and write
$A=\exp(20k^2)$ and $M=1600A$. We construct a path from $a$ to $a'$
with at most
\begin{equation}
 20000A\max\{1,\log(R/D)\}
 \label{eq:input-path-bound}
\end{equation}
steps. Every step is either a nearest-neighbor step or an increment
whose endpoint-set distance from $B$, divided by its length, is
at least $A$. This is the alternative needed for the source's
increment decomposition.

In the disk of radius $R$ centered at $(a+a')/2$, choose $z$ outside
the at most four cones of half-angle $0.1$ starting at $x\in\{a,a'\}$ and
pointing toward $y\in B\setminus\{x\}$. The total area of those cones within the
disk is at most $4(0.1)(3R/2)^2=0.9R^2$, so such a point exists.
The segment from $x$ to $z$ has length $L\le3R/2$. At arclength $s$
along it, elementary triangle geometry gives distance at least
$(d_x+s)/100$ from $B$, where $d_x=\dist(x,B)$: the excluded angle
bounds the distance to each $y\ne x$ by a fixed fraction of $|x-y|+s$.
If $y=x$, the distance is exactly $s$, so the same bound holds.
Round all segment points by one fixed nearest-lattice rule, with error
at most one, using the same rounded endpoint $z$ on both segments.

If $d_x<M$, start with a shortest Manhattan path to the rounded point
at $s_0=\min\{L,M\}$; it takes at most $2M+2$ nearest-neighbor
steps. If $d_x\ge M$, put $s_0=0$. On the remaining segment, increase
$s$ at each step by $(d_x+s)/(800A)$, or stop at its endpoint if that
comes first. With $\rho=d_x+s\ge M$, the rounded step has length
at most $\rho/(400A)$ and both rounded endpoints have distance at
least $\rho/200$ from $B$. Its separation ratio is therefore at least
$2A$. Delete repeated consecutive lattice points.

Geometric growth of $d_x+s$ bounds the number of these sampled steps
on a leg by
\[
 1+1600A\log\frac{d_x+L}{d_x+s_0}.
\]
Whenever that leg is nonempty its denominator is at least $D$.
Concatenating the first leg with the reversed second leg gives at most
\[
 6400A+6+3200A\log\left(1+\frac{3R}{2D}\right)
 \le20000A\max\{1,\log(R/D)\}
\]
steps, proving \cref{eq:input-path-bound}. The case $a=a'$ uses the
empty path. The source orders the two increment lengths; if the other
increment collapses, that ordering reduces to this same case. All
constructed paths are finite and thus lie in sufficiently large
exhausting volumes. Since $S'=\log(D/R)$, this supplies exactly the
positive factor used in Equation~(367), with a constant depending
only on $k$.

\subsection{Slope, free energy, and the Gaussian multiplier}
\label{sec:input-amplitude}

Let $I_{\rm col}$ be the left-minus-right arrow count in one vertical
column of an $M$ by $L$ torus. Conservation makes it independent of
the column. The unit-height slope is
\[
 s=I_{\rm col}/L=I_{\rm total}/(ML),
 \qquad I_{\rm total}=M I_{\rm col}.
\]
The free energy is still normalized per site, by $(ML)^{-1}\log Z$.
Thus the sector in \cite[Definition~4.2]{DKLM2026} is
$I_{\rm col}/L=2\lfloor Ls/2\rfloor/L$. Its printed denominator
$ML$ cannot apply to a \emph{single-column} count: $|I_{\rm col}|\le L$,
so a nonzero fixed slope would then disappear as $M\to\infty$.
The corrected convention agrees with the filling fraction
$(1-s)/2$ in its Equation~(413) and with the per-site sector convention
of \cite[Section~1.2 and Theorem~2]{SixVertexFreeEnergy2022}.

For completeness, the derivative calculation in
\cite[Equations~(412) and~(416)]{DKLM2026} then reads
\[
 f''(0)=-2\left(\frac{\pi-\zeta}{\pi}\right)^2
             \frac{\pi^2}{4(\pi-\zeta)}
       =-\frac{\pi-\zeta}{2}=-\arcsin(c/2).
\]
The endpoint $c=2$, or $\zeta=0$, uses the endpoint calculation
specified there. No factor of the torus length is part of this
second derivative.

The conditional Gaussian comparison uses Dirichlet energy
$\operatorname{Dir}(H)=\int|\nabla H|^2$, without a factor $1/2$.
For the neutral capacitary test $\phi_A$ in
\cite[Section~25]{DKLM2026}, the Gaussian variance is
$\sigma^2/\operatorname{Dir}(H_A)$; hence its logarithmic tail
coefficient is $-\operatorname{Dir}(H_A)/(2\sigma^2)$.
The source's Equation~(484) gives
$\operatorname{Dir}(H_A)\le4\rho+C$, with \(C\) independent of \(\rho\).
In its Equation~(482) this yields
$f''(0)\ge-(1+C/(4\rho))/\sigma^2$; letting \(\rho\to\infty\) gives
$f''(0)\ge-1/\sigma^2$.

In the reverse inequality, retain the factor $(1-\varepsilon)^2$
in its Equation~(485), because the circuit threshold in
Equation~(441) is $(1-\varepsilon)k$. Combined with Equation~(486),
this gives $(1-\varepsilon)^4$ in the bound for $f''(0)$.
Only after the prescribed large-scale limits let
$\varepsilon\downarrow0$; the resulting inequality remains
$f''(0)\le-1/\sigma^2$. The geometric normalization in the two
splitting estimates is consistent: Equations~(503)--(504) use
height $2n(k-24)$ across scale $2nN$, whereas
Equations~(523)--(524) use height $4n(k+4)+8$ across the same scale.
These give slopes $(k-24)/N$ and $2(k+4)/N+o(1)$ respectively,
with the circuit level reparametrization already made in
Equations~(440)--(441).

Other directional slips in that argument are read in the sense of the
events defined there: after Equation~(475) the inclusion is
$C_l\subset C_{l-4}$; at Equation~(514) the ridge event is contained
in the successful exploration event; and the good event at
Equation~(559) has \emph{not too many} arms, as defined in
Equation~(562). The preceding $\eta^{-2}$ bound can be retained
in the combinatorial penalty of Equation~(565), giving the required
constant depending on $(\rho,\varepsilon)$. None of these readings
changes the final order of limits or the multiplier.

Consequently $\sigma^2=1/\arcsin(c/2)$, under the full Gaussian
convergence premise stated in \cref{prop:plane-normalization}.
The source GFF has logarithmic coefficient $1/(2\pi)$, so our
coefficient is $\kappa=\sigma^2/(2\pi)=1/a^2$. This explains both
the slope convention and the conversion to $H=K/2$.

\subsection{The conformal growth estimate}
\label{sec:conformal-growth}

The conformal mapping and coefficient-one Green identities used in
\cref{prop:partition,lem:flat-patch} are given in
\cite[Chapter~6, Sections~1.1 and~5.2]{Ahlfors1979}.
Here is the short derivation of the growth inequality from published
univalent-function inputs, so the numerical constant in
\cref{eq:cut-green-bound} does not depend on an unidentified draft
edition. For a normalized univalent map $f$ of the disk, the second
coefficient obeys $|a_2|\le2$
\cite[Equation~(2.68)]{AstalaIwaniecMartin2009}. Apply this to the
normalized Koebe transform at $z$ to obtain
\[
 \left|(1-|z|^2)\frac{f''(z)}{f'(z)}-2\overline z\right|\le4.
\]
At $z=re^{i\vartheta}$ this bounds the radial derivative of
$\log|f'(re^{i\vartheta})|$ by $(4+2r)/(1-r^2)$. Integration gives
\[
 |f'(re^{i\vartheta})|\le\frac{1+r}{(1-r)^3},
 \qquad |f(re^{i\vartheta})|\le\frac r{(1-r)^2}.
\]
The last inequality follows by integrating along a radius. Together
with the quarter theorem
\cite[Theorem~2.10.5]{AstalaIwaniecMartin2009}, this gives precisely
the conformal-radius bounds used in the flat-patch proof, for every
proper simply connected domain, without boundary smoothness.
